% !TEX root = cap07_standalone.tex
% Capítulo 6. Conclusiones
% Estructura según índice de la Dra. Lárraga.

\chapter{Conclusiones}
\label{ch:conclusiones}

\begin{resumen}
Este capítulo sintetiza los hallazgos de la investigación, delimita sus alcances, identifica las limitaciones del modelo y propone direcciones de trabajo futuro, entre ellas el uso del modelo como apoyo a la planeación. Se presentan las respuestas a las preguntas de investigación, la validación de hipótesis y las aportaciones del trabajo al campo del modelado urbano, a la luz de los resultados de la validación quinquenal múltiple (cinco períodos, 2011--2020).
\end{resumen}

% ============================================================
\section{Principales hallazgos}
% ============================================================

Esta sección sintetiza los resultados centrales de la investigación, presenta las respuestas a las preguntas formuladas en el Capítulo~\ref{ch:introduccion} y evalúa el estado de cada hipótesis a la luz de la validación quinquenal múltiple sobre cinco períodos desplazados.

\subsection{Síntesis de la investigación}

Esta investigación abordó el problema de modelar el crecimiento urbano de Querétaro (1984--2020) mediante la integración de tres metodologías complementarias: \textbf{Weight of Evidence (WoE)} para cuantificar empíricamente la evidencia espacial, \textbf{Autómatas Celulares (AC)} para simular dinámicas espaciales locales y una estrategia de \textbf{calibración sistemática} para el ajuste de parámetros. La validación quinquenal múltiple sobre cinco períodos desplazados (2011--2016, 2012--2017, 2013--2018, 2014--2019 y 2015--2020) sugiere que la metodología WoE-AC, bajo las condiciones de datos y calibración descritas en los capítulos anteriores, presenta varios rasgos consistentes. El \textit{pipeline} LBP+K-Means+SVM clasificó de manera automática las 37 capturas RGB anuales (1984--2020) en mapas binarios urbano/no-urbano, con coherencia interna de $88$--$94$\% respecto de las pseudoetiquetas de K-Means en el conjunto de prueba, no frente a una verdad terreno. Sobre esa base, el WoE cuantificó empíricamente la evidencia espacial a partir de las transiciones históricas observadas en 26 años de datos (1984--2010). El modelo obtuvo un \textbf{FoM promedio de $0{,}317$} sobre cinco períodos de validación desplazados, equivalente a aproximadamente $31{,}7$\% en la escala porcentual usada por \textcite{pontius2008comparing}; interpretado frente al espectro empírico multi-sitio de ese estudio (seis aplicaciones con FoM inferior a $0{,}15$ y una sola por encima de $0{,}50$), el resultado se ubica en una \textbf{banda intermedia} de dicho marco (Capítulo~\ref{ch:estado-del-arte}), utilizando únicamente siete variables espaciales derivadas del mapa binario. En conjunto, reprodujo los patrones de expansión urbana con \textit{Accuracy} promedio de $0{,}722$, Kappa promedio de $0{,}445$ y FoM promedio de $0{,}317$, y capturó el mecanismo de difusión por contagio espacial perimetral de forma coherente con los patrones observados.

\subsection{Respuestas a las preguntas de investigación}

\subsubsection{P1: ¿Métricas comparables a la literatura?}

El FoM promedio de $0{,}317$, el Kappa de $0{,}445$ y el IoU de $0{,}633$ (Tabla~\ref{tab:metrics_all_periods}) se sitúan en la banda intermedia del espectro multi-sitio de \textcite{pontius2008comparing}: por encima del subconjunto con FoM inferior a $0{,}15$ y por debajo del único caso que supera $0{,}50$. El umbral $\theta = 0{,}75$ y el peso de vecindad $\alpha = 0{,}50$ son la configuración usada en las cinco ventanas; ensayos preliminares con umbrales menores produjeron sobre-predicción masiva, con un FoM que no es directamente comparable porque se midió bajo otro protocolo de conteo (Capítulo~\ref{ch:resultados-analisis}).

\subsubsection{P2: ¿Qué variables discriminan?}

El \textit{Information Value} muestra que la \textbf{distancia al área urbana} y la \textbf{densidad de vecindad} (tres radios) concentran $0{,}782$ del peso normalizado ($78{,}2$\%), por encima de fragmentación y gradiente. La lectura es coherente con un mecanismo de difusión por contagio espacial.

\subsubsection{P3: ¿Se sostiene el desempeño en las cinco ventanas?}

El FoM varía entre $0{,}222$ y $0{,}378$ según la dinámica observada. El valor más bajo (2011--2016) coincide con la única ventana con crecimiento neto observado negativo. El comportamiento morfológico (expansión perimetral) se reproduce en las cinco ventanas; las métricas píxel a píxel sí dependen de la señal de cambio. Las ventanas son hold-out respecto de 1984--2010 y se solapan entre sí.

\subsubsection{P4: ¿El flujo corre en código abierto y CPU?}

El pipeline completo se ejecutó en Python (NumPy, SciPy, scikit-learn) sobre capturas de Google Earth en escritorio, sin GPU ni software propietario, en un cómputo del orden de tres horas en CPU estándar. No se ensayó una réplica geográfica en otra ciudad; H3 afirma la reproducibilidad de \emph{este} flujo, no su transferibilidad ya medida.

\subsection{Cumplimiento de los objetivos específicos}

Los cinco objetivos específicos planteados en el Capítulo~\ref{ch:introduccion} se cumplieron del modo que se describe a continuación. El primero (OE1) condujo a la \textbf{serie histórica binaria}: se construyó la serie anual 1984--2020 (37 mapas urbano/no-urbano) a partir de imágenes de archivo de Google Earth, con una cadena de clasificación LBP+K-Means+SVM que no requiere verdad terreno etiquetada, y sin software propietario (Sección~\ref{sec:preprocesamiento}). El segundo (OE2) abordó las \textbf{variables espaciales y el WoE}: se calcularon las siete variables derivadas del mapa binario y se adoptó el WoE como función de transición por su descomposición auditable mediante el \textit{Information Value} (Capítulos~\ref{ch:modelo-crecimiento-urbano} y \ref{ch:resultados-analisis}). El tercero (OE3) corresponde al \textbf{diseño e implementación del WoE-AC}: se implementó el autómata celular que simula la transición no urbano $\rightarrow$ urbano sobre la ZMQ (Capítulo~\ref{ch:modelo-crecimiento-urbano}). El cuarto (OE4) cubrió la \textbf{calibración y el protocolo de validación}: se calibraron umbral y peso de vecindad por búsqueda en cuadrícula y se validó el modelo en cinco ventanas quinquenales desplazadas con FoM, Kappa e IoU (Capítulo~\ref{ch:resultados-analisis}). El quinto (OE5) atendió la \textbf{estabilidad temporal y el posicionamiento}: se analizó la estabilidad del desempeño a lo largo de los cinco períodos y se contrastaron los valores de FoM e IoU con los rangos de la literatura de modelos de caja blanca y de caja negra (Capítulos~\ref{ch:resultados-analisis} y \ref{ch:estado-del-arte}).

\subsection{Cumplimiento de los requerimientos del modelo}

El Capítulo~\ref{ch:introduccion} fijó cuatro requerimientos (Tabla~\ref{tab:requerimientos_modelo_1}) y el Capítulo~\ref{ch:estado-del-arte} los usó para calificar la literatura. Corresponde aplicarlos al propio modelo, con el mismo criterio y sin concesiones.

\textbf{R1, simulación histórica con precisión: cumplido de forma parcial.} El modelo reproduce el patrón de cambio con FoM promedio de $0{,}317$, Kappa de $0{,}445$ e IoU de $0{,}633$ en cinco ventanas desplazadas del período de entrenamiento. La reserva es de magnitud, no de localización: el desacuerdo de cantidad concentra el $79{,}4\%$ del desacuerdo total (Sección~\ref{sec:descomposicion-error}), de modo que el modelo acierta mayormente en dónde crece la ciudad y sobre-predice cuánto.

\textbf{R2, soporte a planeación estratégica: cumplido en su formulación mínima.} El modelo produce mapas espacialmente explícitos que admiten lectura comparativa bajo supuestos declarados. No se calibraron escenarios contrafactuales de política ni se cuantificó su efecto, de modo que esos mapas se leen como escenarios y no como planes calibrados, tal como el propio requerimiento acota.

\textbf{R3, replicabilidad y apertura: cumplido.} El flujo completo se ejecuta en Python con bibliotecas de código abierto sobre CPU estándar, en un cómputo del orden de tres horas, con calibración sistemática documentada y sin software propietario ni hardware especializado.

\textbf{R4, interpretabilidad de la función de transición: cumplido.} Los pesos de la función de transición son valores WoE por bin, legibles como evidencia espacial, y su contribución relativa queda ordenada por el \textit{Information Value} de cada variable (Tabla~\ref{tab:iv_weights}). Ninguna decisión del modelo depende de parámetros opacos.

El balance es que tres de los cuatro requerimientos se cumplen y el primero lo hace con una reserva localizada y medida. Esa reserva, la sobre-predicción de la magnitud, es la que orienta las extensiones de la Sección~\ref{sec:trabajo-futuro}.

\subsection{Validación de hipótesis}

La hipótesis general y sus cuatro particulares, formuladas en el Capítulo~\ref{ch:introduccion}, se contrastan a continuación con los resultados del Capítulo~\ref{ch:resultados-analisis}.

\textbf{Hipótesis general.} Un modelo híbrido WoE-AC, calibrado y validado mediante un protocolo multitemporal sobre datos de libre acceso, permite simular de forma interpretable y reproducible el crecimiento urbano de una ciudad intermedia mexicana con un desempeño comparable al reportado en la literatura. \textbf{Estado: apoyada por la evidencia.} El modelo alcanzó un FoM promedio de $0{,}317$ y un Kappa promedio de $0{,}445$ en las cinco ventanas quinquenales desplazadas, valores en la banda intermedia documentada por \textcite{pontius2008comparing}, con implementación íntegra en Python y datos de acceso libre. La Tabla~\ref{tab:hypothesis_validation} recoge el veredicto de las cuatro hipótesis particulares.

\begin{xltabular}{\textwidth}{@{}l X c X@{}}
\caption{Validación de hipótesis particulares}\label{tab:hypothesis_validation}\\
\toprule
\textbf{Hip.} & \textbf{Enunciado (Cap.~\ref{ch:introduccion})} & \textbf{Estado} & \textbf{Evidencia} \\
\midrule
\endfirsthead
\toprule
\textbf{Hip.} & \textbf{Enunciado (Cap.~\ref{ch:introduccion})} & \textbf{Estado} & \textbf{Evidencia} \\
\midrule
\endhead
\bottomrule
\endlastfoot
H1 & \textit{Interpretabilidad.} La regla de transición WoE ponderada por IV produce un modelo auditable: la distancia al frente urbano y la densidad de vecindad concentran el mayor IV, por encima de fragmentación y gradiente. & Apoyada & IV concentrado en distancia al área urbana y densidad de vecindad ($0{,}782$ del peso normalizado), por encima de fragmentación y gradiente. El orden entre distancia y densidad depende del tratamiento de las celdas no elegibles (Cap.~\ref{ch:resultados-analisis}). \\
H2 & \textit{Estabilidad temporal.} El desempeño se sostiene en las cinco ventanas quinquenales desplazadas; el valor más bajo se asocia a la única ventana con crecimiento neto observado negativo, no a un fallo de calibración. & Parcial & FoM entre $0{,}222$ y $0{,}378$; el mínimo (2011--2016) coincide con crecimiento neto observado negativo. El patrón morfológico se sostiene; las métricas píxel a píxel varían. \\
H3 & \textit{Reproducibilidad.} El flujo completo se ejecuta con herramientas de código abierto en CPU estándar, sin GPU ni software propietario. & Apoyada & Implementación íntegra en Python (NumPy, SciPy, scikit-learn) con capturas de Google Earth en escritorio; cómputo del orden de tres horas en CPU (Sección~\ref{sec:accesibilidad}). \\
H4 & \textit{Desempeño comparable.} Aplicado a la ZMQ, el modelo reproduce el patrón de crecimiento con FoM, Kappa e IoU en la banda intermedia del espectro de \textcite{pontius2008comparing}. & Apoyada & FoM promedio $0{,}317$, Kappa $0{,}445$ e IoU $0{,}633$. El contraste con caja negra es cualitativo (sitios y horizontes distintos). \\
\end{xltabular}

\subsection{Aportaciones del trabajo}

El uso del Weight of Evidence y del autómata celular no constituye, por sí mismo, una novedad metodológica, pues ambos están ampliamente documentados en la literatura. La aportación central de esta investigación reside en el \textbf{protocolo de evaluación}: el acoplamiento WoE-AC \textbf{calibrado y validado de forma reproducible sobre cinco ventanas temporales desplazadas}, con código y datos de acceso libre, aplicado a una ciudad intermedia mexicana poco estudiada. El WoE se incorpora como un campo probabilístico continuo ponderado por el \textit{Information Value}, pero el valor diferenciador está en poder distinguir un desempeño estable de un ajuste a un único período. De esa aportación principal se derivan varias contribuciones específicas. La primera es el \textbf{protocolo de validación quinquenal múltiple}, de carácter \textit{metodológico}: la evaluación sobre cinco períodos desplazados (2011--2020) con métricas píxel a píxel (FoM, Kappa, IoU), la descomposición del error de Pontius entre desacuerdo de cantidad y de asignación y un análisis visual comparativo de patrones morfológicos evalúa la estabilidad temporal del modelo con mayor rigor que la validación en un único período, y es la aportación que sostiene las demás. La segunda, de orden \textit{práctico}, es la \textbf{reproducibilidad abierta del pipeline}: un sistema modular en Python que transforma capturas RGB en mapas binarios y en proyecciones de crecimiento urbano (clasificación híbrida, WoE, AC y evaluación), ejecutable con herramientas de código abierto y datos de acceso libre, documentado y publicado. La tercera, de orden \textit{instrumental} y subordinada al protocolo, es el \textbf{campo WoE continuo ponderado por Information Value}: la implementación concreta que hace auditable la regla de transición, donde la idoneidad WoE no es un mapa estático, sino un campo probabilístico ponderado por el \textit{Information Value} normalizado de cada variable, que se recalcula en cada paso sobre el estado actualizado del mapa y se acopla a la vecindad de Moore. La cuarta, de naturaleza \textit{aplicada}, es la \textbf{aplicación a una ciudad intermedia mexicana}: la adaptación del enfoque WoE-AC al contexto de la Zona Metropolitana de Querétaro, un caso poco estudiado y sin un \textit{benchmark} comparable previo. La quinta, de carácter \textit{empírico}, es la \textbf{serie temporal de Querétaro}: un conjunto de datos procesado de 37 años (1984--2020) de mapas binarios urbano/no-urbano, que puede utilizarse como base para investigaciones futuras en la región. La sexta, igualmente \textit{empírica}, es el \textbf{posicionamiento frente a la literatura}: el contraste del modelo WoE-AC con estudios internacionales mediante métricas estandarizadas (FoM, IoU), que contextualiza los resultados en el espectro empírico multi-sitio de \textcite{pontius2008comparing}.

% ============================================================
\section{Alcances}
% ============================================================

El modelo se aplicó a la Zona Metropolitana de Querétaro sobre una serie anual de 37 años (1984--2020), período que concentra la mayor expansión urbana documentada para esta ciudad y para el cual fue posible construir una serie histórica completa a partir de imágenes de archivo de Google Earth. En su dimensión temporal, el alcance abarca esa serie procesada hasta mapas binarios urbano/no urbano compatibles con el autómata celular. En su dimensión espacial, cada año se representa como una rejilla de $1\,792 \times 3\,024$ píxeles derivada del preprocesamiento de las capturas RGB. El modelo no reemplaza los instrumentos de planeación urbana ni produce predicciones normativas: es una herramienta de análisis retrospectivo y prospectivo con la que se exploran escenarios de crecimiento bajo supuestos explícitos y reproducibles. La metodología fue diseñada para ser replicable en otras ciudades intermedias mexicanas con datos equivalentes, previa recalibración de parámetros.

\subsection{Accesibilidad computacional}
\label{sec:accesibilidad}

El pipeline completo (clasificación + WoE + simulación + validación) se ejecuta en aproximadamente 3 horas en CPU estándar, sin requerir GPU. Los enfoques de aprendizaje profundo en teledetección suelen exigir conjuntos anotados grandes y, en la práctica, hardware gráfico \parencite{Ma2019}; esta tesis no midió horas de GPU de terceros. La Tabla~\ref{tab:approach_characteristics} resume las características operativas del enfoque adoptado.

\begin{table}[htbp]
\centering
\caption{Características operativas del enfoque WoE-AC frente a modelos de aprendizaje profundo}
\label{tab:approach_characteristics}
\begin{tabularx}{\textwidth}{@{}>{\raggedright\arraybackslash}p{0.28\textwidth} X X@{}}
\toprule
\textbf{Característica} & \textbf{WoE-AC (este trabajo)} & \textbf{Aprendizaje profundo} \\
\midrule
Datos de entrenamiento requeridos & Series multitemporales locales (en este trabajo, 26 transiciones anuales agregadas por \textit{pooling}). & Conjuntos de datos etiquetados a gran escala, habitualmente con decenas de imágenes anotadas. \\
Hardware mínimo & CPU estándar. & GPU recomendada. \\
Interpretabilidad de resultados & Alta (pesos WoE explicables). & Baja (modelo de caja negra). \\
Tiempo de ejecución & Del orden de tres horas en CPU. & Depende del conjunto y de la arquitectura; suele requerir GPU. \\
Transferibilidad & Requiere recalibración WoE. & Requiere reentrenamiento. \\
\bottomrule
\end{tabularx}
\end{table}

% ============================================================
\section{Limitaciones del modelo}
% ============================================================

El rigor académico exige reconocer los límites del enfoque adoptado. Esta sección describe las limitaciones conceptuales, empíricas y de transferibilidad del modelo, e identifica las extensiones metodológicas que permitirían superarlas en investigaciones futuras.

\subsection{Limitaciones de alcance del estudio}

El modelo opera sobre imágenes RGB sin bandas infrarrojas, lo que restringe la clasificación binaria a información espectral visible. Esta decisión fue deliberada: reducir la barrera de entrada a datos de acceso libre y procesamiento reproducible, a costa de la precisión espectral que ofrece la teledetección multiespectral. La representación binaria urbano/no urbano simplifica la complejidad interna de la ciudad; como consecuencia, el modelo captura la dinámica de expansión de la mancha urbana, pero no su estructura interna.

El enfoque es puramente espacial: factores como políticas públicas, incentivos económicos o decisiones de inversión influyeron en el crecimiento real de Querétaro y no están representados en las variables de entrada. Esta limitación es compartida por la mayoría de los modelos CA-WoE de la literatura y acota la interpretación de los resultados al patrón morfológico observado, no a sus causas socioeconómicas.

La elección de caja blanca implica una concesión de desempeño frente a modelos de aprendizaje profundo, cuya magnitud este trabajo no midió. Los resultados del Capítulo~\ref{ch:resultados-analisis} ubican al modelo WoE-AC en la banda intermedia del espectro multi-sitio de \textcite{pontius2008comparing}; el contraste con modelos de caja negra es cualitativo, porque se refiere a sitios, clases y horizontes distintos, de modo que no permite cuantificar el costo de la interpretabilidad.

\subsection{Limitaciones conceptuales}

En el plano conceptual, la primera limitación atañe a la \textbf{definición operativa de ``urbano''}. El clasificador K-Means+SVM, al operar sobre imágenes RGB sin banda infrarroja, aprende a separar zonas con cubierta vegetal verde de zonas sin ella. La clase ``urbano'' incluye, por tanto, tejido urbano real pero también suelo desnudo agrícola en temporada seca y superficies áridas. Una clasificación más precisa requeriría imágenes con banda NIR para calcular NDVI real, o muestras de campo geolocalizadas como ground truth. Esta limitación no invalida la validación interna del modelo (que evalúa la consistencia predictiva dentro de la misma clasificación), pero sí acota la interpretación de los resultados en términos de área urbana absoluta. Una segunda limitación es el \textbf{radio de vecindad}: Moore ($3 \times 3$) captura interacciones locales pero no influencias de largo alcance (empleo, transporte). Los \emph{pesos} WoE, estimados en 1984--2010, quedan fijos; las siete variables espaciales sí se recalculan en cada paso sobre el estado actualizado del mapa. Por último, la \textbf{asunción de estacionariedad}: las relaciones WoE calibradas con 1984--2010 se asumen válidas en las cinco ventanas posteriores, supuesto que puede no sostenerse ante transformaciones económicas o políticas disruptivas.

\subsection{Limitaciones empíricas}

En el plano empírico, la primera limitación se relaciona con la \textbf{calidad de los datos}: la clasificación sin verdad terreno (K-Means + SVM sobre pseudoetiquetas) tiene una \textit{accuracy} interna del $88$ al $94$\%, lo que equivale a un error nominal del $6$ al $12$\% por imagen (del orden de $325\,000$ a $650\,000$ píxeles sobre $5{,}4 \times 10^{6}$). Estos errores se propagan al cálculo WoE y pueden amplificarse en la simulación del AC mediante el efecto de vecindad. La segunda es la \textbf{resolución espectral}: las imágenes contienen únicamente bandas RGB, sin información infrarroja, y el índice verde--rojo normalizado (NGRDI), $(G-R)/(G+R)$, es un proxy cromático menos discriminativo que el NDVI real $(NIR-R)/(NIR+R)$. La tercera concierne a las \textbf{variables internas}: el modelo usa 7 variables espaciales derivadas del propio mapa binario (densidades, distancias, fragmentación, gradiente) y no incorpora variables externas potencialmente relevantes como distancia a vialidades, pendiente topográfica, zonificación oficial o variables socioeconómicas. La cuarta es el \textbf{crecimiento disperso no modelado}: el mecanismo de difusión por vecindad no captura el desarrollo tipo leapfrog (saltos discontinuos de urbanización), que representa una fracción significativa del crecimiento real observado.

\subsection{Limitaciones de transferibilidad}

En cuanto a la transferibilidad, una primera limitación es la \textbf{especificidad contextual}: los pesos WoE son específicos al patrón histórico de Querétaro y requieren recalibración para su aplicación en otras ciudades. La segunda es la \textbf{dependencia de la resolución}: los resultados están ligados a la grilla de los mapas binarios derivada de las capturas homogeneizadas ($1\,792 \times 3\,024$ píxeles), de modo que modificar la extensión, el tamaño de celda o la fuente de imágenes implica reentrenar el WoE y recalibrar los parámetros del AC. La tercera es el \textbf{horizonte temporal}: la aplicabilidad del modelo para horizontes mayores a 15 años es incierta sin recalibración periódica, dado que las dinámicas urbanas pueden cambiar estructuralmente.

% ============================================================
\section{Trabajo futuro}
\label{sec:trabajo-futuro}
% ============================================================

Con base en las limitaciones identificadas, se proponen las siguientes extensiones prioritarias:

\subsection{Corto plazo (1-2 años)}

A corto plazo destacan tres líneas. La primera es la incorporación de \textbf{variables adicionales}: distancia a vialidades principales, distancia a centros de empleo y pendiente topográfica mediante análisis WoE multi-variable. La segunda es la \textbf{optimización metaheurística}, esto es, explorar metaheurísticas alternativas al grid exhaustivo (p.\,ej.\ optimización por enjambre de partículas o procedimientos bayesianos de búsqueda) para la calibración de parámetros, comparando su desempeño con el procedimiento empleado en esta investigación. La tercera es la \textbf{validación en más ciudades}: aplicar la metodología a Aguascalientes, León y San Luis Potosí para evaluar la transferibilidad geográfica.

\subsection{Mediano plazo (3-5 años)}

A mediano plazo se perfilan cuatro extensiones. El \textbf{WoE temporal} consistiría en desarrollar pesos que evolucionan dinámicamente para capturar cambios en las dinámicas de crecimiento a lo largo del tiempo. La \textbf{integración multi-escala} acoplaría el modelo AC con modelos económicos regionales que proporcionen demanda agregada de urbanización. Los \textbf{datos de movilidad} permitirían incorporar patrones de origen-destino de viajes (disponibles vía telefonía móvil) como variable explicativa. La \textbf{calibración multiobjetivo}, por su parte, extendería el grid search para optimizar simultáneamente FoM, similitud morfológica y compacidad urbana como criterios de sostenibilidad.

\subsection{Largo plazo (5-10 años)}

A largo plazo, la \textbf{integración con cambio climático} adaptaría el modelo para incorporar escenarios de cambio climático (temperatura, precipitación) como restricciones o drivers de urbanización. El \textbf{transfer learning} buscaría desarrollar pre-entrenamiento en múltiples ciudades para reducir los datos requeridos en nuevas aplicaciones. Finalmente, las \textbf{métricas complementarias de validación} pasarían por adoptar de forma sistemática métricas espaciales adicionales (FRAGSTATS, índice de Moran, \textit{Fuzzy Accuracy}) para complementar las métricas píxel a píxel ya empleadas, en línea con buenas prácticas reconocidas en la literatura LUCC.

\subsection{Refinamientos del estimador WoE}

Dos ajustes del procedimiento de estimación son de implementación inmediata. Ninguno corrige un error de los resultados reportados, porque los pesos afectados no intervienen en la simulación; ambos depuran la interpretación de los \textit{Information Value}. El primero es restringir el muestreo a las celdas elegibles, es decir, a las no urbanas en el año inicial de cada par: con ello desaparecen los \textit{bins} sin transiciones observadas, los IV vuelven a la escala convencional de interpretación y la jerarquía entre variables queda libre del efecto trivial que hoy la condiciona (Sección~\ref{ssec:alcance-iv}). El segundo es emparejar cada configuración espacial anual con su transición correspondiente, en lugar de promediar las variables sobre el período de entrenamiento. Ese cambio mostraría además si los pesos son estables entre décadas o si el régimen espacial de la zona cambió a lo largo de los 26 períodos.

\subsection{Uso en planeación urbana}

Cualquier uso instrumental exige antes recalibración y validación local. Cumplida esa condición, las líneas de aplicación son acotadas: anticipar dónde extender redes de agua, drenaje y electricidad; reservar suelo para espacio público y amortiguamiento ambiental antes de que la presión inmobiliaria lo encarezca; y detectar áreas ecológicamente sensibles situadas en la trayectoria de expansión proyectada. En los tres casos el producto es una lectura de tendencia, no un plan.

La estructura del modelo también admite explorar escenarios de política ajustando parámetros o aplicando máscaras de probabilidad, por ejemplo densificación intraurbana, contención del borde o concentración a lo largo de corredores de transporte. Conviene ser explícito en el alcance: este trabajo no calibró corridas contrafactuales por política ni cuantificó su efecto sobre métricas de dispersión para Querétaro. Ese ejercicio queda como extensión, y su primer requisito es definir cómo se traduce cada instrumento de política en un parámetro del autómata.

\section{Reflexión de cierre}

El aporte concreto de este trabajo es un protocolo de validación quinquenal múltiple, reproducible sobre datos de acceso libre, ejercitado en un modelo WoE-AC de la Zona Metropolitana de Querétaro construido sobre 37 años de imágenes de archivo (1984--2020). Ese protocolo, más que un valor de desempeño aislado, es lo que distingue la propuesta: separa un comportamiento espacial estable de un ajuste a un único período, en un contexto, las ciudades intermedias mexicanas, donde rara vez se documenta esa distinción.

Un modelo parsimonioso e interpretable, validado empíricamente en varios períodos desplazados, puede reproducir de forma razonable la dinámica morfológica del crecimiento urbano y servir como herramienta de exploración. Con siete variables derivadas del propio mapa binario, el FoM promedio de $0{,}317$ sobre cinco ventanas quinquenales desplazadas se ubica en una banda intermedia del espectro multi-sitio documentado en la literatura \parencite{pontius2008comparing}. El mecanismo de difusión por contagio espacial se reproduce con razonable fidelidad sin requerir variables socioeconómicas ni hardware especializado.

La metodología no es un fin en sí misma, sino una base reproducible: código abierto, datos libres y validación multitemporal. Sobre esa base, trabajos futuros pueden mejorar la clasificación, incorporar variables socioeconómicas y extender la evaluación a otras ciudades intermedias mexicanas. Su utilidad para la planificación es indirecta: ofrece una primera lectura de hacia dónde tiende a expandirse la ciudad, no una predicción operativa, que exigiría recalibración local y validación adicional antes de cualquier uso instrumental.
